% Original problems 95--103 in section 2.3 of user PDF; printed pp.73--74.
\thermokineticproblem{95}{73}
利用麦克斯韦分布计算气体分子
速率倒数的平均值
$\overline{1/v}$，温度为 $T$、分子质量为 $m$。
将其与平均速率的倒数 $1/\bar v$ 比较。
\exsolution
速率概率密度为
\[
F(v)=4\pi\left(\frac{m}{2\pi k_BT}\right)^{3/2}
v^2e^{-mv^2/(2k_BT)},\qquad v\ge0.
\]
因此
\[
\overline{\frac1v}=\int_0^\infty
\frac1vF(v)\,\mathrm dv
=4\pi\left(\frac m{2\pi k_BT}\right)^{3/2}
\int_0^\infty ve^{-mv^2/(2k_BT)}\,\mathrm dv.
\]
用
$\int_0^\infty ve^{-av^2}\mathrm dv=1/(2a)$，
得到
\[
\boxed{\overline{\frac1v}=
\sqrt{\frac{2m}{\pi k_BT}}}.
\]
又
$\bar v=\sqrt{8k_BT/(\pi m)}$，
所以
\[
\boxed{\overline{\frac1v}
=\frac4\pi\,\frac1{\bar v}}.
\]
\excheck 由于函数 $1/v$ 是凸函数，
Jensen 不等式要求
$\overline{1/v}\ge1/\bar v$，
而 $4/\pi>1$，方向正确。

\thermokineticproblem{96}{73}
单分子质量为 $m$、温度为 $T$ 的理想气体
服从麦克斯韦速率分布。
求平动动能 $\varepsilon=mv^2/2$ 的概率分布密度；
求最概然动能 $\varepsilon_p$，
并判断它是否等于最概然速率对应的动能。
\exsolution
速率概率密度为
$F(v)=4\pi(m/(2\pi k_BT))^{3/2}v^2
e^{-mv^2/(2k_BT)}$。
变量变换
$\varepsilon=mv^2/2$，
$v=\sqrt{2\varepsilon/m}$，
$\mathrm dv/\mathrm d\varepsilon=1/(mv)$，
所以
\[
f_\varepsilon(\varepsilon)=
F(v)\left|\frac{\mathrm dv}{\mathrm d\varepsilon}\right|
=\boxed{\frac2{\sqrt\pi}
\frac{\sqrt\varepsilon}{(k_BT)^{3/2}}
e^{-\varepsilon/(k_BT)}},
\qquad\varepsilon\ge0.
\]
对其对数求导：
\[
\frac{\mathrm d}{\mathrm d\varepsilon}\ln f_\varepsilon
=\frac1{2\varepsilon}-\frac1{k_BT}=0,
\]
因此
\[
\boxed{\varepsilon_p=\frac12k_BT}.
\]
而速率分布的众数为
$v_p=\sqrt{2k_BT/m}$，
对应动能
$mv_p^2/2=k_BT$，并不等于 $\varepsilon_p$。
\excheck 众数在非线性变量变换下不保持不变；
原因是概率密度还要乘雅可比因子。

\thermokineticproblem{97}{73}
求热平衡理想单原子气体中，
单分子平动动能与平均动能之差
不超过 $\delta=1.0\%$
平均动能的分子比例。
\exsolution
三维平动平均动能为
$\bar\varepsilon=3k_BT/2$，
动能分布为第96题的
\[
f_\varepsilon(\varepsilon)
=\frac{2\sqrt\varepsilon}{\sqrt\pi(k_BT)^{3/2}}
e^{-\varepsilon/(k_BT)}.
\]
目标区间
$[(1-\delta)\bar\varepsilon,
(1+\delta)\bar\varepsilon]$。
因 $\delta\ll1$，
把概率密度视作区间中央的常值，
得到
\[
P\simeq2\delta\bar\varepsilon\,
f_\varepsilon(\bar\varepsilon).
\]
代入
$\bar\varepsilon=3k_BT/2$：
\[
\boxed{P\simeq
\frac{6\delta}{\sqrt\pi}
\sqrt{\frac32}e^{-3/2}
\simeq0.0092\simeq0.9\%}.
\]
\excheck 概率小于约2%的相对能量窗口，
因为单分子能量分布较宽，
不能把分子全都视为接近平均动能。

\thermokineticproblem{98}{73}
在温度为 $T$ 的理想气体中，
求单分子平动动能
$\varepsilon>\varepsilon_0$ 的比例，
特别研究 $\varepsilon_0\gg k_BT$ 的情形。
\exsolution
由麦克斯韦能量分布
\[
f_\varepsilon(\varepsilon)
=\frac{2\sqrt\varepsilon}{\sqrt\pi(k_BT)^{3/2}}
e^{-\varepsilon/(k_BT)},
\]
所求尾部概率
\[
P=\frac2{\sqrt\pi}\int_{z_0}^\infty
z^{1/2}e^{-z}\,\mathrm dz,\qquad
z_0=\frac{\varepsilon_0}{k_BT}.
\]
利用不完全伽马函数恒等式，
精确写成
\[
P=\operatorname{erfc}\sqrt{z_0}
+\frac2{\sqrt\pi}\sqrt{z_0}e^{-z_0}.
\]
当 $z_0\gg1$，主导项为
\[
\boxed{P\simeq\frac2{\sqrt\pi}
\sqrt{\frac{\varepsilon_0}{k_BT}}
e^{-\varepsilon_0/(k_BT)}}.
\]
\excheck 高能尾部虽有
$\sqrt{\varepsilon_0/(k_BT)}$ 因子，
但最终由指数衰减主导。

\thermokineticproblem{99}{73}
分子从容器小孔形成的射流，
其速率密度满足
$F(v)=A v^3e^{-mv^2/(2k_BT)}$，
其中 $A$ 为归一化常数。
求射流分子的
(a) 最概然速率及其与容器内最概然速率之比；
(b) 最概然平动动能。
\exsolution
(a) 在 $v>0$ 时
\[
\frac{\mathrm d}{\mathrm dv}\ln F(v)
=\frac3v-\frac{mv}{k_BT}.
\]
令其为零：
\[
v_{p,\rm beam}^2=\frac{3k_BT}{m},
\quad
\boxed{v_{p,\rm beam}=\sqrt{\frac{3k_BT}{m}}}.
\]
容器内气体分子的最概然速率为
$v_{p,\rm gas}=\sqrt{2k_BT/m}$，
故
\[
\boxed{\frac{v_{p,\rm beam}}{v_{p,\rm gas}}
=\sqrt{\frac32}}.
\]
(b) 令 $\varepsilon=mv^2/2$，
概率密度变量变换含
$\mathrm dv/\mathrm d\varepsilon=1/(mv)$，
所以
$f_\varepsilon(\varepsilon)\propto
v^2e^{-\varepsilon/(k_BT)}
\propto\varepsilon e^{-\varepsilon/(k_BT)}$。
取对数求导，
$1/\varepsilon-1/(k_BT)=0$，得
\[
\boxed{\varepsilon_{p,\rm beam}=k_BT}.
\]
\excheck 射流速率分布额外包含入射速率
通量加权，故比容器内的 Maxwell 速率分布
多一因子 $v$。

\thermokineticproblem{100}{73}
温度为 $T$、分子质量为 $m$、
数密度为 $n$ 的理想气体，
求单位墙面积每单位时间以
与墙法线夹角在
$[\theta,\theta+\mathrm d\theta]$
之间的方向撞击墙的分子数。
\exsolution
在气体平衡状态，
速度方向均匀分布在立体角上，
其密度为 $\mathrm d\Omega/(4\pi)$。
取墙法线为极轴，入射仅来自朝向墙的半球，
方向立体角
$\mathrm d\Omega=2\pi\sin\theta\,\mathrm d\theta$。
碰撞通量再乘法向速率
$v\cos\theta$，对速率平均后得
\[
\mathrm d\Phi=
n\bar v\cos\theta\,
\frac{2\pi\sin\theta\,\mathrm d\theta}{4\pi}
=\frac12n\bar v\sin\theta\cos\theta\,\mathrm d\theta.
\]
因为
$\bar v=\sqrt{8k_BT/(\pi m)}$，
\[
\boxed{\mathrm d\Phi=
n\sqrt{\frac{2k_BT}{\pi m}}\,
\sin\theta\cos\theta\,\mathrm d\theta},
\quad 0\le\theta\le\frac\pi2.
\]
\excheck 对半球入射角区间积分，
$\int_0^{\pi/2}\sin\theta\cos\theta\,\mathrm d\theta=1/2$，
恢复总通量 $n\bar v/4$。

\thermokineticproblem{101}{74}
沿用第100题的理想气体与壁面模型，
求单位墙面积每单位时间
以速率在 $[v,v+\mathrm dv]$ 的区间内
撞击墙的分子数。
\exsolution
气体分子速率密度
\[
F(v)=4\pi\left(\frac m{2\pi k_BT}\right)^{3/2}
v^2e^{-mv^2/(2k_BT)}.
\]
对于固定速率 $v$，
入射半球方向的法向速度加权平均为
$v/4$；因此对应碰撞通量
\[
\mathrm d\Phi=
n\frac v4F(v)\,\mathrm dv.
\]
代入分布函数：
\[
\boxed{\mathrm d\Phi=
n\pi\left(\frac m{2\pi k_BT}\right)^{3/2}
v^3e^{-mv^2/(2k_BT)}\,\mathrm dv}.
\]
\excheck 必须有额外通量权重 $v$；
其速率分布与第99题的射流分布同型，
而非容器内部的 $v^2$ 型。

\thermokineticproblem{102}{74}
一种粒子在均匀保守力场中，
沿力的反方向相隔
$\Delta h=3.0\,\mathrm{cm}$ 的两位置，
粒子浓度相差
$\eta=2.0$ 倍。
体系温度 $T=280\,\mathrm K$，
求每个粒子所受外力大小。
\exsolution
玻尔兹曼分布
$n(h)\propto e^{-U(h)/(k_BT)}$。
如果均匀力大小为 $F$，
沿反力方向升高 $\Delta h$
会使势能增加
$\Delta U=F\Delta h$。
故浓度比
\[
\eta=\frac{n_{\rm low}}{n_{\rm high}}
=e^{F\Delta h/(k_BT)}.
\]
取对数：
\[
\boxed{F=\frac{k_BT\ln\eta}{\Delta h}}.
\]
代入
$k_B=1.380649\times10^{-23}\,\mathrm{J/K}$、
$T=280\,\mathrm K$、$\Delta h=0.030\,\mathrm m$，
\[
\boxed{F\simeq8.9\times10^{-20}\,\mathrm N}.
\]
\excheck 由浓度\emph{比值}求力，
不必知道浓度的绝对归一化值。

\thermokineticproblem{103}{74}
显微镜观察悬浮的藤黄胶体微滴，
在相隔 $h=40\,\mathrm{\mu m}$
的两水平层中，
平均微滴数密度相差 $\eta=2.0$ 倍。
环境温度 $T=290\,\mathrm K$。
微滴直径 $d=0.40\,\mathrm{\mu m}$，
其密度比周围液体大
$\Delta\rho=0.20\,\mathrm{g/cm^3}$。
由这些数据求阿伏伽德罗常数 $N_A$。
\exsolution
一个球形微滴的体积
$V_d=\pi d^3/6$。
考虑浮力后的\emph{有效重力}
\[
F_{\rm eff}=\Delta\rho\,V_d g
=\frac{\pi d^3}{6}\Delta\rho\,g.
\]
根据玻尔兹曼分布，
\[
\ln\eta=\frac{F_{\rm eff}h}{k_BT}
=\frac{\pi d^3\Delta\rho gh}
{6k_BT}.
\]
由 $R=N_Ak_B$，
\[
\boxed{N_A=\frac{6RT\ln\eta}
{\pi d^3\Delta\rho gh}}.
\]
代入
$d=4.0\times10^{-7}\,\mathrm m$、
$\Delta\rho=200\,\mathrm{kg/m^3}$、
$h=4.0\times10^{-5}\,\mathrm m$，
计算得
\[
\boxed{N_A\simeq6.4\times10^{23}\,\mathrm{mol^{-1}}}.
\]
\excheck 这是经典佩兰胶体沉降方法的思路；
若漏掉浮力所对应的液体密度，
结果将系统性偏离。
